\section{Explicit redundancy-nine slice cover}\label{app:r9-slices}

This appendix prints the finite polynomial data used in
Proposition~\ref{prop:r9-components}.  All coefficients lie in $\F_7$.
For a vector $[c_0,\ldots,c_m]$, write
\[
 [c_0,\ldots,c_m](L)=\sum_{j=0}^m c_jL^j.
\]
On the squarefree quartic atlas $h_L=[1,0,L,0,1]$, the following
ordered four-point bases give the residual branch discriminants
$\Delta_i(L)$.

\begingroup
\scriptsize
\setlength{\tabcolsep}{3pt}
\begin{longtable}{c p{0.14\textwidth} p{0.65\textwidth}}
\toprule
$i$ & fixed roots & coefficients of $\Delta_i$, constant term first\\
\midrule
\endfirsthead
\toprule
$i$ & fixed roots & coefficients of $\Delta_i$, constant term first\\
\midrule
\endhead
\bottomrule
\endfoot
1 & \texttt{[0, 1, 2, 3]} &
\ttfamily[6, 0, 6, 0, 3, 4, 5, 5, 0, 0, 2, 6, 1, 3, 3, 0, 0, 4, 5, 2, 6, 0, 0, 1, 1]\\
2 & \texttt{[0, 1, 2, 4]} &
\ttfamily[0, 0, 4, 5, 1, 5, 2, 3, 1, 6, 1, 5, 4, 6, 0, 3, 0, 0, 4, 1, 2, 0, 6, 0, 4]\\
3 & \texttt{[0, 1, 2, 5]} &
\ttfamily[2, 6, 2, 1, 1, 3, 0, 1, 6, 5, 0, 2, 3, 4, 4, 4, 5, 4, 0, 4, 6, 3, 1, 5, 5]\\
4 & \texttt{[0, 1, 2, 6]} &
\ttfamily[0, 0, 0, 0, 0, 0, 3, 0, 0, 6, 5, 0, 0, 1, 0, 0, 2, 4, 0, 0, 3, 0, 0, 6, 5]\\
5 & \texttt{[0, 1, 3, 4]} &
\ttfamily[6, 1, 3, 4, 6, 5, 4, 6, 6, 4, 0, 0, 2, 6, 1, 0, 1, 2, 5, 4, 1, 5, 4, 3, 5]\\
6 & \texttt{[1, 2, 3, 4]} &
\ttfamily[6, 5, 5, 6, 0, 4, 0, 5, 6, 3, 6, 6, 0, 3, 4, 1, 4, 2, 3]\\
\end{longtable}

The B\'ezout coefficients in \eqref{eq:r9-bezout} are:
\begin{longtable}{c p{0.84\textwidth}}
\toprule
$i$ & coefficients of $b_i$, constant term first\\
\midrule
\endfirsthead
\toprule
$i$ & coefficients of $b_i$, constant term first\\
\midrule
\endhead
\bottomrule
\endfoot
1 & \ttfamily[6, 4, 1, 2, 3, 5, 5, 1, 6, 5, 0, 0, 4, 5, 3, 6, 2, 1, 1, 4, 3, 1, 2, 6, 4, 2, 6, 4, 6, 3, 0, 5, 5, 2, 6]\\
2 & \ttfamily[1, 0, 5, 1, 1, 5, 4, 2, 6, 5, 0, 4, 6, 4, 4, 4, 1, 3, 0, 1, 3, 3, 5, 1, 6, 4, 3, 6, 5, 0, 2, 4, 4, 5, 2]\\
3 & \ttfamily[]\\
4 & \ttfamily[]\\
5 & \ttfamily[]\\
6 & \ttfamily[0, 3, 5, 4, 0, 1]\\
\end{longtable}
\endgroup

Only $b_1,b_2,b_6$ are nonzero, but all six sections are retained as
regression controls.  Direct multiplication in $\F_7[L]$ gives
\eqref{eq:r9-bezout}, so the principal opens $D(\Delta_i)$ cover the
squarefree atlas.

For completeness, the multiple-root boundary has the following four
normal forms.  The branch vector is again written constant term first.
\[
\begin{array}{c|c|c|c|c}
\text{partition}&h&\text{fixed roots}&\text{reduced branch}&
\text{discriminant}\\ \hline
4 &[1,0,0,0,0]&[1,2,3,4]&[2,2,1]&3\\
3+1 &[0,1,0,0,0]&[0,2,4,6]&[4,4,3]&3\\
2+2 &[0,0,1,0,0]&[0,1,2,3]&[5,5,2,1]&1\\
2+1+1 &[0,1,6,0,0]&[0,1,2,3]&[3,5,6,6,6]&5
\end{array}
\]
Every reduced discriminant is nonzero.  The squarefree atlas and
these four boundary strata therefore exhaust all nonzero geometric
binary quartics used in the surjectivity argument.
